\documentclass[11pt]{article}
\usepackage[utf8]{inputenc}
\usepackage[T1]{fontenc}
\usepackage{amsmath,amssymb,mathtools}
\usepackage[margin=1in]{geometry}
\usepackage{hyperref}
\usepackage{array}
\usepackage{parskip}
\setlength{\parindent}{0pt}
\hypersetup{colorlinks=true,linkcolor=blue,urlcolor=blue}
\title{Lecture 17 — Inverse Laplace Transform and Properties}
\author{}
\date{}
\begin{document}
\maketitle

\textbf{Course:} CEC 315 — Signals and Systems (Spring 2026)
\textbf{Instructor:} Rogelio Gracia Otalvaro
\textbf{Source PDF:} \texttt{all\_lectures/cec315-lctr17-inverse-laplace-properties.pdf}
\textbf{Exam coverage:} Exam 3

\textbf{Lctr 17: Inverse Laplace Transform, Pole-Zero Analysis, and Properties}

Rogelio Gracia Otalvaro
Spring 2026

\textbf{Lctr 17: Inverse Laplace Transform, Pole-Zero Analysis, and Properties}

Spring 2026

Focus: Sections 9.3–9.6 (Pages 670–692)

\section*{17.1 The Big Picture}

\begin{quote}
\textbf{Why This Matters}

In Lecture 16 we learned how to go \textit{forward}: given a time-domain signal $x(t)$, compute its Laplace transform $X(s)$ and ROC. In practice, we almost always need to go in the \textit{reverse} direction: given a system's transfer function $H(s)$ (which we get from a differential equation), find the corresponding impulse response $h(t)$.

This lecture covers three essential skills: (1) how to invert the Laplace transform using partial fractions, (2) how to read system behavior directly from pole-zero plots without doing algebra, and (3) the properties of the Laplace transform that let us handle operations like differentiation, convolution, and time shifting algebraically.
\end{quote}

\textbf{Roadmap:}

\begin{enumerate}
  \item The inverse Laplace transform via partial fraction expansion.
  \item Worked examples: distinct real poles, repeated poles, complex poles.
  \item Geometric evaluation of the Fourier transform from pole-zero plots.
  \item Properties of the Laplace transform (with ROC implications).
  \item Summary table of transform pairs and properties.
\end{enumerate}

\section*{17.2 The Inverse Laplace Transform}

\subsection*{17.2.1 The Formal Inversion Integral}

The exact inverse is:

\[x(t) = \frac{1}{2\pi j} \int_{\sigma-j\infty}^{\sigma+j\infty} X(s)\, e^{st}\, ds\]

This is a contour integral in the complex plane along a vertical line in the ROC. In this course, we will \textbf{never evaluate this integral directly}. Instead, we use a far more practical method.

\subsection*{17.2.2 The Practical Method: Partial Fraction Expansion}

The strategy is:

\begin{enumerate}
  \item Express $X(s)$ as a sum of simple terms $\frac{A_i}{s-p_i}$, $\frac{A_i}{(s-p_i)^n}$, etc.
  \item Invert each term using the table from Lecture 16.
  \item Use the ROC to determine whether each term is right-sided or left-sided.
\end{enumerate}

\begin{quote}
\textbf{Key Insight}

The ROC tells you which direction each exponential goes. For a pole at $s = p_i$:

- If the ROC is to the \textbf{right} of $p_i$: the term is $e^{p_i t} u(t)$ (right-sided).
- If the ROC is to the \textbf{left} of $p_i$: the term is $-e^{p_i t} u(-t)$ (left-sided).

This is the only place where the ROC enters the inversion process, but it is critical.
\end{quote}

\section*{17.3 Example: Distinct Real Poles}

\begin{quote}
\textbf{Inverse Laplace Transform with Distinct Real Poles}

Find $x(t)$ given:

\[X(s) = \frac{2s+6}{(s+1)(s+3)}, \qquad \text{Re}\{s\} > -1\]

\textbf{Step 1: Partial fraction expansion.}

\[\frac{2s+6}{(s+1)(s+3)} = \frac{A}{s+1} + \frac{B}{s+3}\]

Multiply both sides by $(s+1)(s+3)$:

\[2s+6 = A(s+3) + B(s+1)\]

\textbf{Cover-up method:}

- Set $s = -1$: $2(-1)+6 = A(-1+3) \Rightarrow 4 = 2A \Rightarrow A = 2$.
- Set $s = -3$: $2(-3)+6 = B(-3+1) \Rightarrow 0 = -2B \Rightarrow B = 0$.

So: $X(s) = \dfrac{2}{s+1} + \dfrac{0}{s+3} = \dfrac{2}{s+1}$

\textbf{Step 2: Check the ROC for each pole.}
The ROC is $\text{Re}\{s\} > -1$. The pole at $s = -1$ has the ROC to its \textit{right}, so the corresponding term is right-sided.

\textbf{Step 3: Invert.}

\[\boxed{x(t) = 2\, e^{-t}\, u(t)}\]

\textbf{Check:} At $t = 0^+$: $x(0^+) = 2$. From the initial value theorem (which we will see later), $\lim_{s \to \infty} s X(s) = \lim_{s \to \infty} \frac{2s^2 + 6s}{s^2 + 4s + 3} = 2$. \checkmark
\end{quote}

\section*{17.4 Example: Distinct Real Poles with Mixed ROC}

\begin{quote}
\textbf{Two-Sided Signal from Partial Fractions}

Find $x(t)$ given:

\[X(s) = \frac{5s+17}{(s+1)(s-3)}, \qquad -1 < \text{Re}\{s\} < 3\]

\textbf{Step 1: Partial fractions.}

\[\frac{5s+17}{(s+1)(s-3)} = \frac{A}{s+1} + \frac{B}{s-3}\]

- $s = -1$: $5(-1)+17 = A(-1-3) \Rightarrow 12 = -4A \Rightarrow A = -3$.
- $s = 3$: $5(3)+17 = B(3+1) \Rightarrow 32 = 4B \Rightarrow B = 8$.

So: $X(s) = \dfrac{-3}{s+1} + \dfrac{8}{s-3}$

\textbf{Step 2: Determine the direction of each term from the ROC.}
The ROC is the strip $-1 < \sigma < 3$.

- Pole at $s = -1$: The ROC is to the \textbf{right} of this pole $\Rightarrow$ right-sided.
- Pole at $s = 3$: The ROC is to the \textbf{left} of this pole $\Rightarrow$ left-sided.

\textbf{Step 3: Invert each term.}

\[\frac{-3}{s+1}\ (\text{right-sided}) \longrightarrow -3\, e^{-t}\, u(t)\]

\[\frac{8}{s-3}\ (\text{left-sided}) \longrightarrow -8\, e^{3t}\, u(-t)\]

\[\boxed{x(t) = -3\, e^{-t}\, u(t) - 8\, e^{3t}\, u(-t)}\]
\end{quote}

\textbf{Figure:} Pole-zero plot in the $s$-plane showing poles at $s = -1$ (marked with $\times$) and $s = 3$ (marked with $\times$). The ROC is the vertical strip between $\sigma = -1$ and $\sigma = 3$ (shaded green), bounded by dashed vertical lines. Labels below the strip indicate "right-sided" (for the pole at $-1$) and "left-sided" (for the pole at $3$).

\section*{17.5 Example: Repeated Poles}

\begin{quote}
\textbf{Repeated Poles}

Find $x(t)$ given:

\[X(s) = \frac{4s+5}{(s+1)^2 (s+3)}, \qquad \text{Re}\{s\} > -1\]

\textbf{Step 1: Partial fraction expansion.}
For a repeated pole at $s = -1$ of order 2, we need:

\[\frac{4s+5}{(s+1)^2(s+3)} = \frac{A}{s+1} + \frac{B}{(s+1)^2} + \frac{C}{s+3}\]

Multiply through by $(s+1)^2 (s+3)$:

\[4s+5 = A(s+1)(s+3) + B(s+3) + C(s+1)^2\]

\textbf{Find $C$:} set $s = -3$: $4(-3)+5 = C(-3+1)^2 = 4C \Rightarrow C = -7/4$.
\textbf{Find $B$:} set $s = -1$: $4(-1)+5 = B(-1+3) = 2B \Rightarrow B = 1/2$.
\textbf{Find $A$:} compare coefficients of $s^2$: $0 = A + C \Rightarrow A = -C = 7/4$.

So:

\[X(s) = \frac{7/4}{s+1} + \frac{1/2}{(s+1)^2} + \frac{-7/4}{s+3}\]

\textbf{Step 2:} ROC is $\text{Re}\{s\} > -1$. All poles are at $\sigma \le -1$, so all terms are right-sided.

\textbf{Step 3: Invert} using the pair $\dfrac{1}{(s+a)^2} \leftrightarrow t\, e^{-at}\, u(t)$:

\[\boxed{x(t) = \frac{7}{4} e^{-t} u(t) + \frac{1}{2} t\, e^{-t} u(t) - \frac{7}{4} e^{-3t} u(t)}\]
\end{quote}

\begin{quote}
\textbf{Warning}

For repeated poles of order $n$, the partial fraction expansion includes terms $\frac{B_1}{(s-p)}, \frac{B_2}{(s-p)^2}, \ldots, \frac{B_n}{(s-p)^n}$. The inverse of $\frac{1}{(s-p)^k}$ involves $t^{k-1} e^{pt}$, not just $e^{pt}$. Forgetting the lower-order terms $\frac{B_1}{s-p}, \ldots$ is a very common mistake.
\end{quote}

\section*{17.6 Example: Complex Conjugate Poles}

\begin{quote}
\textbf{Complex Poles: Second-Order System}

Find $x(t)$ given:

\[X(s) = \frac{2s}{s^2 + 4s + 13}, \qquad \text{Re}\{s\} > -2\]

\textbf{Step 1: Find the poles} by completing the square:

\[s^2 + 4s + 13 = (s+2)^2 + 9 = (s+2)^2 + 3^2\]

Poles are at $s = -2 \pm j3$ (complex conjugate pair with $\sigma_0 = -2$, $\omega_d = 3$).

\textbf{Step 2: Rewrite the numerator} in terms of $(s+2)$:

\[2s = 2(s+2) - 4\]

So:

\[X(s) = \frac{2(s+2)}{(s+2)^2 + 9} - \frac{4}{(s+2)^2 + 9} = 2 \cdot \frac{(s+2)}{(s+2)^2 + 3^2} - \frac{4}{3} \cdot \frac{3}{(s+2)^2 + 3^2}\]

\textbf{Step 3: Invert} using the pairs:

\[\frac{s+a}{(s+a)^2 + \omega_d^2} \ \overset{\mathcal{L}}{\longleftrightarrow}\ e^{-at} \cos(\omega_d t)\, u(t)\]

\[\frac{\omega_d}{(s+a)^2 + \omega_d^2} \ \overset{\mathcal{L}}{\longleftrightarrow}\ e^{-at} \sin(\omega_d t)\, u(t)\]

\[\boxed{x(t) = \left[2\, e^{-2t} \cos(3t) - \frac{4}{3}\, e^{-2t} \sin(3t)\right] u(t)}\]

This is a damped sinusoid: the $e^{-2t}$ envelope decays (since the real part of the poles is $-2 < 0$), and the oscillation frequency is $\omega_d = 3$ rad/s.
\end{quote}

\textbf{Figure:} Plot of $x(t) = e^{-2t}\left[2\cos(3t) - \frac{4}{3}\sin(3t)\right] u(t)$ versus $t$. The signal starts at $x(0^+) = 2$, oscillates, and decays toward zero. A red dashed envelope $\pm A e^{-2t}$ bounds the damped oscillation. Horizontal $t$-axis from 0 to 4, vertical axis from $-1$ to $2$.

\begin{quote}
\textbf{Key Insight}

Complex conjugate poles always come in pairs $s = -\sigma_0 \pm j\omega_d$. They produce \textbf{damped sinusoidal} time-domain signals of the form $e^{-\sigma_0 t}\left[A\cos(\omega_d t) + B\sin(\omega_d t)\right] u(t)$. The real part $-\sigma_0$ controls the decay rate; the imaginary part $\omega_d$ controls the oscillation frequency. This is the same behavior we saw in second-order systems in Lecture 15, now derived systematically from the Laplace transform.
\end{quote}

\section*{17.7 Geometric Evaluation from Pole-Zero Plots}

\subsection*{17.7.1 The Key Formula}

Given $X(s)$ in factored form:

\[X(s) = K \frac{\prod_{i=1}^{M} (s - z_i)}{\prod_{i=1}^{N} (s - p_i)}\]

Each factor $(s - z_i)$ or $(s - p_i)$ is a complex number. We can write it in polar form:

\[(s - z_i) = |s - z_i|\, e^{j \angle (s - z_i)}\]

The magnitude and phase of $X(s)$ are therefore:

\[\boxed{|X(s)| = |K| \frac{\prod_{i=1}^{M} |s - z_i|}{\prod_{i=1}^{N} |s - p_i|}} \qquad \boxed{\angle X(s) = \angle K + \sum_{i=1}^{M} \angle (s - z_i) - \sum_{i=1}^{N} \angle (s - p_i)}\]

Geometrically, $|s - p_i|$ is the \textbf{distance} from the point $s$ to the pole $p_i$ in the $s$-plane, and $\angle (s - p_i)$ is the \textbf{angle} of the vector from $p_i$ to $s$.

\subsection*{17.7.2 Evaluating the Frequency Response}

To get the frequency response $H(j\omega)$, we evaluate at $s = j\omega$ (points on the $j\omega$-axis):

\textbf{Figure:} $s$-plane diagram showing the $j\omega$-axis (vertical) and $\sigma$-axis (horizontal). A zero ($\circ$) labeled $z_1$ sits near the origin and a complex conjugate pair of poles ($\times$) labeled $p_1$ and $p_1^*$ are in the left half plane. A test point at $s = j\omega$ is marked on the $j\omega$-axis. Three vectors are drawn: from $z_1$ to $s = j\omega$ (labeled $|s - z_1|$), from $p_1$ to $s = j\omega$ (labeled $|s - p_1|$), and from $p_1^*$ to $s = j\omega$ (labeled $|s - p_1^*|$). The formula annotated is $|H(j\omega)| = |K| \dfrac{|s - z_1|}{|s - p_1| \cdot |s - p_1^*|}$.

\begin{quote}
\textbf{Key Insight}

\textbf{Near a pole:} the denominator vector gets short $\Rightarrow |H|$ gets large (resonance peak).
\textbf{Near a zero:} the numerator vector gets short $\Rightarrow |H|$ gets small (notch/null).
This gives you a quick qualitative picture of the frequency response by just looking at the pole-zero plot.
\end{quote}

\subsection*{17.7.3 First-Order System Example}

Consider $H(s) = \dfrac{1}{s+a}$ with $a > 0$ (single pole at $s = -a$, no finite zeros):

\[|H(j\omega)| = \frac{1}{|j\omega + a|} = \frac{1}{\sqrt{a^2 + \omega^2}}\]

\textbf{Figure:} Plot of $|H(j\omega)| = 1/\sqrt{a^2 + \omega^2}$ for $a = 1, 2, 4$. Three curves are shown: $a = 1$ (solid blue), $a = 2$ (dashed red), and $a = 4$ (dotted green). All curves peak on the $j\omega$-axis at $\omega = 0$, with peak values $1$, $0.5$, and $0.25$ respectively. Horizontal axis $\omega$ ranges from $-10$ to $10$.

When the pole is farther from the $j\omega$-axis (larger $a$), the minimum distance from any point on the $j\omega$-axis to the pole is larger, so the peak magnitude is lower and the response is wider. This connects directly to the bandwidth–decay-rate relationship from Lecture 15.

\section*{17.8 Properties of the Laplace Transform}

In the table below, $x(t) \overset{\mathcal{L}}{\longleftrightarrow} X(s)$ with ROC $R_x$, and $y(t) \overset{\mathcal{L}}{\longleftrightarrow} Y(s)$ with ROC $R_y$.

\begin{center}
\begin{tabular}{|l|l|l|}
\hline
\textbf{Property} & \textbf{Time Domain} & \textbf{$s$-Domain / ROC} \\
\hline
Linearity & $ax(t) + by(t)$ & $aX(s) + bY(s)$, ROC $\supseteq R_x \cap R_y$ \\
\hline
Time shifting & $x(t - t_0)$ & $e^{-st_0} X(s)$, ROC $= R_x$ \\
\hline
$s$-domain shifting & $e^{s_0 t} x(t)$ & $X(s - s_0)$, ROC $= R_x$ shifted by $\text{Re}\{s_0\}$ \\
\hline
Time scaling & $x(at)$, $a > 0$ & $\dfrac{1}{a} X\!\left(\dfrac{s}{a}\right)$, ROC $= a R_x$ \\
\hline
Conjugation & $x^*(t)$ & $X^*(s^*)$, ROC $= R_x$ \\
\hline
Convolution & $x(t) * y(t)$ & $X(s) \cdot Y(s)$, ROC $\supseteq R_x \cap R_y$ \\
\hline
Differentiation in $t$ & $\dfrac{d}{dt} x(t)$ & $s X(s)$, ROC $\supseteq R_x$ \\
\hline
Integration in $t$ & $\displaystyle\int_{-\infty}^{t} x(\tau)\, d\tau$ & $\dfrac{1}{s} X(s)$, ROC $\supseteq R_x \cap \{\text{Re}\{s\} > 0\}$ \\
\hline
Differentiation in $s$ & $-t\, x(t)$ & $\dfrac{d}{ds} X(s)$, ROC $= R_x$ \\
\hline
\end{tabular}
\end{center}

\subsection*{17.8.1 Why These Properties Matter}

\begin{quote}
\textbf{Key Insight}

\textbf{Convolution $\longleftrightarrow$ Multiplication:} This is the most important property for system analysis. If a system has impulse response $h(t) \leftrightarrow H(s)$ and input $x(t) \leftrightarrow X(s)$, then the output is:

\[y(t) = x(t) * h(t) \ \overset{\mathcal{L}}{\longleftrightarrow}\ Y(s) = X(s) \cdot H(s)\]

Convolution in time becomes multiplication in $s$, the same simplification as the Fourier transform, but now applicable to a much wider class of signals.

\textbf{Differentiation $\longleftrightarrow$ Multiplication by $s$:} This is why the Laplace transform turns differential equations into algebraic equations. Each derivative $d/dt$ becomes a multiplication by $s$.
\end{quote}

\subsection*{17.8.2 Worked Example: Using the $s$-Domain Shift Property}

\begin{quote}
\textbf{Using $s$-Domain Shifting}

Find the Laplace transform of $x(t) = e^{-3t} \cos(5t)\, u(t)$.

\textbf{Step 1:} Start with the known pair:

\[\cos(5t)\, u(t)\ \overset{\mathcal{L}}{\longleftrightarrow}\ \frac{s}{s^2 + 25}, \qquad \text{Re}\{s\} > 0\]

\textbf{Step 2:} Apply the $s$-domain shift: multiplication by $e^{-3t}$ in time $\Rightarrow$ replace $s$ by $s + 3$ in the transform:

\[e^{-3t} \cos(5t)\, u(t)\ \overset{\mathcal{L}}{\longleftrightarrow}\ \frac{s+3}{(s+3)^2 + 25}, \qquad \text{Re}\{s\} > -3\]

Notice the ROC shifted left by 3: from $\sigma > 0$ to $\sigma > -3$.
\end{quote}

\subsection*{17.8.3 Worked Example: Using the Differentiation Property}

\begin{quote}
\textbf{Transform of $t\, e^{-2t} u(t)$ via the Differentiation-in-$s$ Property}

\textbf{Step 1:} Start with:

\[e^{-2t} u(t)\ \overset{\mathcal{L}}{\longleftrightarrow}\ \frac{1}{s+2}, \qquad \text{Re}\{s\} > -2\]

\textbf{Step 2:} Apply the property $-t\, x(t) \leftrightarrow \dfrac{d}{ds} X(s)$:

\[-t\, e^{-2t} u(t)\ \overset{\mathcal{L}}{\longleftrightarrow}\ \frac{d}{ds}\!\left(\frac{1}{s+2}\right) = \frac{-1}{(s+2)^2}\]

\textbf{Step 3:} Multiply both sides by $-1$:

\[\boxed{t\, e^{-2t} u(t)\ \overset{\mathcal{L}}{\longleftrightarrow}\ \frac{1}{(s+2)^2}, \qquad \text{Re}\{s\} > -2}\]

This confirms the repeated-pole pair from the table.
\end{quote}

\section*{17.9 The Initial-Value and Final-Value Theorems}

These theorems let you extract the behavior of $x(t)$ at $t = 0^+$ and $t \to \infty$ directly from $X(s)$, without inverting.

\begin{quote}
\textbf{Initial-Value Theorem}

If $x(t) = 0$ for $t < 0$ and $x(t)$ contains no impulses at $t = 0$:

\[\boxed{x(0^+) = \lim_{s \to \infty} s\, X(s)}\]
\end{quote}

\begin{quote}
\textbf{Final-Value Theorem}

If $x(t) = 0$ for $t < 0$ and $x(t)$ has a finite limit as $t \to \infty$ (i.e., all poles of $s X(s)$ have negative real parts):

\[\boxed{x(\infty) = \lim_{s \to 0} s\, X(s)}\]
\end{quote}

\begin{quote}
\textbf{Warning}

The final-value theorem is only valid if $\lim_{t \to \infty} x(t)$ actually exists as a finite number. If $x(t)$ oscillates (poles on the $j\omega$-axis) or grows (poles in the RHP), the theorem gives a \textbf{meaningless result}. Always check that $s X(s)$ has all its poles in the open left half-plane before applying it.
\end{quote}

\begin{quote}
\textbf{Initial and Final Value Theorems}

Let $X(s) = \dfrac{10}{s(s+2)(s+5)}$, ROC: $\text{Re}\{s\} > 0$.

\textbf{Initial value:}

\[x(0^+) = \lim_{s \to \infty} s \cdot \frac{10}{s(s+2)(s+5)} = \lim_{s \to \infty} \frac{10}{(s+2)(s+5)} = 0\]

\textbf{Final value:} Check: $sX(s) = \dfrac{10}{(s+2)(s+5)}$ has poles at $s = -2, -5$ (both in LHP). \checkmark

\[x(\infty) = \lim_{s \to 0} \frac{10}{(s+2)(s+5)} = \frac{10}{(2)(5)} = 1\]

So $x(t)$ starts at 0 and settles to 1 as $t \to \infty$, a step response that rises from zero to a final value of 1.
\end{quote}

\section*{17.10 Comprehensive Laplace Transform Pairs}

\begin{center}
\begin{tabular}{|l|l|l|}
\hline
\textbf{Signal $x(t)$} & \textbf{$X(s)$} & \textbf{ROC} \\
\hline
$\delta(t)$ & $1$ & all $s$ \\
\hline
$u(t)$ & $\dfrac{1}{s}$ & $\text{Re}\{s\} > 0$ \\
\hline
$t\, u(t)$ & $\dfrac{1}{s^2}$ & $\text{Re}\{s\} > 0$ \\
\hline
$\dfrac{t^{n-1}}{(n-1)!}\, u(t)$ & $\dfrac{1}{s^n}$ & $\text{Re}\{s\} > 0$ \\
\hline
$e^{-at}\, u(t)$ & $\dfrac{1}{s+a}$ & $\text{Re}\{s\} > -\text{Re}\{a\}$ \\
\hline
$t\, e^{-at}\, u(t)$ & $\dfrac{1}{(s+a)^2}$ & $\text{Re}\{s\} > -\text{Re}\{a\}$ \\
\hline
$e^{-at} \cos(\omega_d t)\, u(t)$ & $\dfrac{s+a}{(s+a)^2 + \omega_d^2}$ & $\text{Re}\{s\} > -a$ \\
\hline
$e^{-at} \sin(\omega_d t)\, u(t)$ & $\dfrac{\omega_d}{(s+a)^2 + \omega_d^2}$ & $\text{Re}\{s\} > -a$ \\
\hline
\end{tabular}
\end{center}

\section*{17.11 Summary}

\begin{enumerate}
  \item \textbf{Partial fractions} is the practical method for inverting Laplace transforms: decompose $X(s)$ into simple terms, invert each using the table, and use the ROC to determine right-sided vs. left-sided.
  \item \textbf{Distinct real poles} give pure exponentials. \textbf{Repeated poles} give polynomial-times-exponential terms. \textbf{Complex poles} give damped sinusoids.
  \item The \textbf{pole-zero plot} gives geometric insight: the frequency response magnitude is the product of zero-distances divided by the product of pole-distances. Near a pole $\Rightarrow$ large magnitude (resonance). Near a zero $\Rightarrow$ small magnitude (notch).
  \item The \textbf{convolution property} $(x * h \leftrightarrow XH)$ is the foundation of system analysis. The \textbf{differentiation property} $(dx/dt \leftrightarrow sX)$ turns differential equations into algebra.
  \item The \textbf{initial-} and \textbf{final-value theorems} let you check boundary behavior without full inversion.
\end{enumerate}

\section*{17.12 Common Mistakes to Avoid}

\begin{enumerate}
  \item \textbf{Forgetting to use the ROC to assign direction:} Each term in the partial fraction expansion can be right-sided or left-sided. The ROC, not the sign of the pole, determines which.
  \item \textbf{Incomplete partial fractions for repeated poles:} A pole of order $n$ at $s = p$ requires $n$ terms: $\dfrac{B_1}{s-p}, \dfrac{B_2}{(s-p)^2}, \ldots, \dfrac{B_n}{(s-p)^n}$.
  \item \textbf{Wrong handling of complex poles:} Never split complex conjugate poles into separate $\dfrac{A}{s-p}$ and $\dfrac{A^*}{s-p^*}$ terms and try to invert them individually with real exponentials. Instead, complete the square and use the cos/sin pairs.
  \item \textbf{Applying the final-value theorem when it does not apply:} If $sX(s)$ has poles on the $j\omega$-axis or in the RHP, the theorem is invalid. Always check first.
  \item \textbf{Confusing pole-zero distances with pole-zero locations:} The frequency response depends on \textit{distances from $s = j\omega$ to each pole/zero}, not on the pole/zero locations themselves.
\end{enumerate}

Rogelio Gracia Otalvaro

\section*{Practice Problems Summary}

\begin{enumerate}
  \item \textbf{Example (17.3) — Distinct Real Poles:} Find $x(t)$ given $X(s) = \dfrac{2s+6}{(s+1)(s+3)}$, $\text{Re}\{s\} > -1$. Answer: $x(t) = 2 e^{-t} u(t)$.
  \item \textbf{Example (17.4) — Distinct Real Poles with Mixed ROC:} Find $x(t)$ given $X(s) = \dfrac{5s+17}{(s+1)(s-3)}$, $-1 < \text{Re}\{s\} < 3$. Answer: $x(t) = -3 e^{-t} u(t) - 8 e^{3t} u(-t)$.
  \item \textbf{Example (17.5) — Repeated Poles:} Find $x(t)$ given $X(s) = \dfrac{4s+5}{(s+1)^2 (s+3)}$, $\text{Re}\{s\} > -1$. Answer: $x(t) = \frac{7}{4} e^{-t} u(t) + \frac{1}{2} t e^{-t} u(t) - \frac{7}{4} e^{-3t} u(t)$.
  \item \textbf{Example (17.6) — Complex Conjugate Poles:} Find $x(t)$ given $X(s) = \dfrac{2s}{s^2 + 4s + 13}$, $\text{Re}\{s\} > -2$. Answer: $x(t) = \left[2 e^{-2t} \cos(3t) - \frac{4}{3} e^{-2t} \sin(3t)\right] u(t)$.
  \item \textbf{Example (17.7.3) — First-Order Frequency Response:} Evaluate $|H(j\omega)|$ for $H(s) = \dfrac{1}{s+a}$, $a > 0$. Answer: $|H(j\omega)| = 1/\sqrt{a^2 + \omega^2}$.
  \item \textbf{Example (17.8.2) — $s$-Domain Shifting:} Find the Laplace transform of $x(t) = e^{-3t} \cos(5t) u(t)$. Answer: $\dfrac{s+3}{(s+3)^2 + 25}$, $\text{Re}\{s\} > -3$.
  \item \textbf{Example (17.8.3) — Differentiation in $s$:} Derive the Laplace transform of $t\, e^{-2t} u(t)$ using the differentiation-in-$s$ property. Answer: $\dfrac{1}{(s+2)^2}$, $\text{Re}\{s\} > -2$.
  \item \textbf{Example (17.9) — Initial and Final Value Theorems:} For $X(s) = \dfrac{10}{s(s+2)(s+5)}$, $\text{Re}\{s\} > 0$, find $x(0^+)$ and $x(\infty)$. Answer: $x(0^+) = 0$, $x(\infty) = 1$.
\end{enumerate}

\end{document}
